\documentclass[11pt]{article}
\usepackage{fontspec}
\usepackage[french]{babel}
\usepackage{amsmath}
\usepackage{amssymb}
\usepackage{amsthm}
\usepackage{mathrsfs}
\usepackage{tikz-cd}
\usepackage{enumitem}

\newcommand{\printedpage}[1]{\marginpar{\footnotesize p.~#1}}
\newcommand{\msp}[1]{\marginpar{\footnotesize [#1]}}
\newcommand{\cX}{\mathscr{X}}
\newcommand{\Ob}{\operatorname{Ob}}

\begin{document}

\printedpage{150}
% head: A. Grothendieck : Dérivateurs, Ch. XII, Édition des Universités de Montpellier II et Paris VII
% foot: Version 30 mars 2003

Reste à prouver c). Je vais le déduire d'une

\textbf{Proposition 17.} \textit{Soit un carré commutatif}
\[
Q \quad = \quad
\begin{tikzcd}
X_0 \arrow[r, "i_1"] \arrow[d, "i_2"'] & X_1 \arrow[d, "f_1"] \\
X_2 \arrow[r, "f_2"] & Y
\end{tikzcd}
\]
\textit{dans} Cat\textit{, et $\cX_Q$ la catégorie fibrée sur}
\[
\tilde{\Phi} \quad (\simeq \mathbf{\Delta}^1 \times \mathbf{\Delta}^1) \quad = \quad
\begin{tikzcd}
0 & 1 \arrow[l] \\
2 \arrow[u] & 3 \arrow[l] \arrow[u]
\end{tikzcd}
\]
\textit{qui l'intègre. Soit $\cX$ l'image inverse du fermé}
\[
\begin{tikzcd}
0 & 1 \arrow[l] \\
2 \arrow[u] &
\end{tikzcd}
\]
\msp{page 235}
\textit{dans $\cX_Q$, qui intègre le diagramme}
\[
\begin{tikzcd}[row sep=small]
 & X_1 \\
X_0 \arrow[ru] \arrow[rd] & \\
 & X_2
\end{tikzcd}
\]
\textit{. Pour que $Q$ soit $W$-cocartésien, il faut et il suffit que l'inclusion}
\[
i : \cX \longrightarrow \cX_Q
\]
\textit{soit dans $W$.}

Comme $\cX_Q$ est lisse sur $\tilde{\Phi}$, et 3 est objet initial de $\tilde{\Phi}$, il s'ensuit que l'inclusion de la fibre de $\cX_Q$ en le sommet 3 de $\tilde{\Phi}$ est dans $W$,
\[
Y \xrightarrow{i_Y} \cX_Q \qquad\qquad \in W.
\]

Considérons alors le diagramme
\[
\begin{tikzcd}[column sep=small]
\cX \arrow[rr, "\varphi_Q"] \arrow[rd, "i"'] & & Y \arrow[ld, "i_Y\in W"] \arrow[ll, dotted, shift left=3ex, shorten <=1.2em, shorten >=1.2em, "\alpha"'] \\
 & \cX_Q. &
\end{tikzcd}
\]

\printedpage{151}
% head: A. Grothendieck : Dérivateurs, Ch. XII, Édition des Universités de Montpellier II et Paris VII
% foot: Version 30 mars 2003

Ce diagramme n'est pas commutatif, mais il y a une flèche de commutativité
\[
\alpha : i_Y\varphi_Q \longrightarrow i.
\]

On sait que ceci implique que
\[
i \in W \qquad\qquad \Longleftrightarrow \qquad\qquad (i_Y\varphi_Q) \in W
\]
(cf. lemme p. 2), et comme $i_Y \in W$,
\[
(i_Y\varphi_Q) \in W \qquad\qquad \Longleftrightarrow \qquad\qquad \varphi_Q \in W,
\]
d'où
\[
i \in W \qquad\qquad \Longleftrightarrow \qquad\qquad \varphi_Q \in W,
\]

\hfill q.e.d.

Nous pouvons maintenant prouver

\msp{page 236}

\textbf{Proposition 18} (\footnote[127]{devrait être vrai dans toute catégorie de modèles ayant assez de cofibrations.}). (Moyennant W(7 bis).) \textit{Considérons un diagramme commutatif}
\[
\begin{tikzcd}
X_0 \arrow[r] \arrow[d] \arrow[rd, phantom, "1"] & X_1 \arrow[r] \arrow[d] \arrow[rd, phantom, "2"] & X_2 \arrow[d] \\
X'_0 \arrow[r] & X'_1 \arrow[r] & X'_2
\end{tikzcd}
\tag*{$(\ast)$}
\]
\textit{d'où un carré composé}
\[
\begin{tikzcd}
X_0 \arrow[r] \arrow[d] \arrow[rd, phantom, "3"] & X_2 \arrow[d] \\
X'_0 \arrow[r] & X'_2.
\end{tikzcd}
\]
\textit{Supposons que le carré} 1 \textit{soit $W$-cocartésien. Alors le carré} 2 \textit{est $W$-cocartésien si et seulement si le carré} 3 \textit{l'est.}

\textsc{Démonstration.}

\texttt{[Maintenant, }$X_i$\texttt{ et }$X'_i$\texttt{ échangent leurs rôles.]}

\printedpage{152}
% head: A. Grothendieck : Dérivateurs, Ch. XII, Édition des Universités de Montpellier II et Paris VII
% foot: Version 30 mars 2003

Soit
\[
I \quad = \quad
\left(
\begin{tikzcd}
0 \arrow[d] & 1 \arrow[l] \arrow[d] & 2 \arrow[l] \arrow[d] \\
0' & 1' \arrow[l] & 2' \arrow[l]
\end{tikzcd}
\right)
\]
et soit $\cX$ l'intégrale du diagramme $(\ast)$, catégorie fibrée sur $\cX$. On désigne par $\cX_{010'}$ la sous-catégorie pleine de $\cX$ image inverse de la sous-catégorie pleine % [sic] fibrée sur $\cX$ -> fibrée sur $I$, a category fibred over itself; the base is the category $I$ just defined
\[
\underbrace{I_{010'}}_{=\{0,1,0'\}} \ \subseteq I,
\]
et de même pour toute autre partie $A$ de $\Ob I = \{0,1,2,0',1',2'\}$, donnant une sous-catégorie pleine $\cX_A$ de $\cX$. On pose
\[
U \quad = \quad
\begin{tikzcd}
0 \arrow[d] & 1 \arrow[l] & 2 \arrow[l] \\
0' & &
\end{tikzcd}
, \qquad
V \quad = \quad
\begin{tikzcd}
0 \arrow[d] & 1 \arrow[l] \arrow[d] & 2 \arrow[l] \\
0' & 1' \arrow[l] &
\end{tikzcd}
,
\]
d'où
\[
\cX_U = \cX_{0120'}, \qquad\qquad \cX_V = \cX_{0120'1'}.
\]
Les flèches entre les $\cX_A$ sont les inclusions

\msp{page 237}

déduites d'inclusions $A \hookrightarrow A'$ dans $\mathfrak{P}(\Ob I)$. Ainsi, le caractère $W$-cocartésien des carrés 1, 2, 3 s'exprime respectivement par la condition que les inclusions
\[
\begin{aligned}
\cX_{010'} \quad &\xrightarrow{\varphi_1} \quad \cX_{010'1'} \\
\cX_{121'} \quad &\xrightarrow{\varphi_2} \quad \cX_{121'2'} \\
\cX_{020'} \quad &\xrightarrow{\varphi_1} \quad \cX_{020'2'} % [sic] $\varphi_1$ -> $\varphi_3$, this inclusion is called $\varphi_3$ in the diagram that follows
\end{aligned}
\]
soient dans $W$. Considérons alors le diagramme de flèches d'inclusion
\printedpage{153}
% head: A. Grothendieck : Dérivateurs, Ch. XII, Édition des Universités de Montpellier II et Paris VII
% foot: Version 30 mars 2003
\[
\begin{tikzcd}
\cX_{020'} \arrow[r, "\varphi_3"] \arrow[d, "\alpha"'] & \cX_{020'2'} \arrow[d, "\beta"] \\
\cX_U \arrow[r, "\psi_3"] \arrow[d, "\psi_1"'] & \cX \arrow[d, "\mathrm{id}"] \\
\cX_V \arrow[r, "\psi_2"] & \cX \\
\cX_{121'} \arrow[r, "\varphi_2"] \arrow[u, "\alpha'"] & \cX_{121'2'}. \arrow[u, "\beta'"']
\end{tikzcd}
\]

Je dis que dans ce diagramme on a
\[
\alpha, \beta, \alpha', \beta' \in W
\tag*{$(\ast)$}
\]
et
\[
\varphi_1 \in W \ \ \text{(i.e. carré 1 $W$-cocartésien)} \qquad\qquad \Longleftrightarrow \qquad\qquad \psi_1 \in W.
\tag*{$(\ast\ast)$}
\]

Les relations $(\ast)$ impliquent déja que % [sic] déja -> déjà, missing grave accent
\[
\begin{aligned}
\varphi_2 \in W \quad &\Longleftrightarrow \quad \psi_2 \in W \\
\varphi_3 \in W \quad &\Longleftrightarrow \quad \psi_3 \in W,
\end{aligned}
\]
et sous l'hypothèse $W$-cocartésienne $(\ast\ast)$, impliquant $\psi_1 \in W$,

\msp{page 238}

on trouve
\[
\psi_2 \in W \qquad\qquad \Longleftrightarrow \qquad\qquad \psi_3 \in W,
\]
d'où
\[
\varphi_2 \in W \qquad\qquad \Longleftrightarrow \qquad\qquad \varphi_3 \in W,
\]
ce qu'il fallait prouver.

Il reste à établir $(\ast)$ et $(\ast\ast)$. Les relations $(\ast)$ proviennent de la remarque que si $i : A \longrightarrow A'$ est une inclusion de sous-catégories pleines de $I$, telle que $i$ soit $W$-\textit{asphérique} (i.e. $A/j$ $W$-asphérique pour $j \in \Ob A'\backslash\Ob A$), alors $\cX_A \longrightarrow \cX_{A'}$ l'est aussi (changement de base lisse), donc est \texttt{[dans]} $W$. On doit vérifier la $W$-asphéricité pour les quatre inclusions
\[
\begin{array}{rcll}
\{0,2,0'\} & \longrightarrow & U = \{0,1,2,0'\} & \qquad\qquad \mathbf{\Delta}^0 \\[1ex]
\{0,2,0',2'\} & \longrightarrow & I & \qquad\qquad \mathbf{\Delta}^0/\mathbf{\Delta}^1 \\[1ex]
\{1,2,1'\} & \longrightarrow & V = \{0,1,2,0',1'\} & \qquad\qquad \mathbf{\Delta}^1,\ \Phi \\[1ex]
\{1,2,1',2'\} & \longrightarrow & I & \qquad\qquad \mathbf{\Delta}^1,\ \mathbf{\Delta}^1 \times \mathbf{\Delta}^1.
\end{array}
\]

\printedpage{154}
% head: A. Grothendieck : Dérivateurs, Ch. XII, Édition des Universités de Montpellier II et Paris VII
% foot: Version 30 mars 2003

(NB J'ai marqué en marge les catégories qu'on obtient comme $A/j$, à isomorphisme près, dans les quatre cas envisagés - elles ont toutes un objet initial, donc sont $W$-asphériques.)

Reste à prouver $(\ast\ast)$. Soit
\[
U' = \{0,1,0',1'\}, \qquad\qquad \text{donc } V = U \cup U',\ \ U \cap U' = \{0,1,0'\},
\]
où $U$, $U'$ sont des \textit{fermés} de $V$ qui recouvrent $V$. Donc on aura de même
\[
\cX_V = \cX_U \cup \cX_{U'}, \qquad\qquad \cX_U \cap \cX_{U'} = \cX_{010'},
\]

\msp{page 239}

\[
\begin{tikzcd}
\cX_{U\cap U'} \arrow[r] \arrow[d, "\varphi_1"] & \cX_U \arrow[d, "\psi_1"] \\
\cX_{U'} \arrow[r] & \cX_V.
\end{tikzcd}
\]

Donc par W(7 bis) (cas dual des immersions fermées), on a
\[
\varphi_1 \in W \qquad\qquad \Longleftrightarrow \qquad\qquad \psi_1 \in W,
\]

\hfill q.e.d.

Notons pour mémoire le corollaire qui nous avait motivé.

\textbf{Corollaire.} \textit{Supposons} W(7 bis) (comme dans toute cette section). \textit{Soit un carré cocartésien dans} Cat
\[
\begin{tikzcd}
X_0 \arrow[r, "i_1"] \arrow[d, "i_2"'] & X_1 \arrow[d] \\
X_2 \arrow[r] & X_3,
\end{tikzcd}
\]
\textit{avec $i_1$ (ou $i_2$) une $W$-cofibration (p.ex. une immersion de Dwyer). Alors ce carré est $W$-cocartésien.}

Cette digression (p. 232 ff.) n'est pas nécessaire d'ailleurs pour formuler dès maintenant la conclusion de nos cogitations sur les cofibres $W$-homotopiques : \textit{Il y a une telle théorie dans} Cat, pourvu seulement que $W$ satisfasse à W(7 bis). En effet, tout résulte du fait que toute flèche $f : X \longrightarrow Y$ dans Cat se factorise en
\[
X \xrightarrow{i_f} \tilde{Y} \xrightarrow{p_f} Y,
\]
avec $i_f$ une $W$-cofibration (et même une imm-

\end{document}
